\chapter{Tangent measures of rectifiable pieces}\label{sec:rectifiable-tangent-transfer}
We prove the forward implication of the tangent-plane characterization in \cite{Sim18}, Chapter 3, Theorem 1.6, printed pp.~87--88. The main issue is that other pieces of a rectifiable set can accumulate near the chosen piece. Differentiation with respect to the measure of their union controls this contribution without assuming a positive distance between pieces. Here $1\leq m\leq n$.

\begin{lemma}[Borel representatives for locally finite Hausdorff measure]\label{lem:rectifiable-borel-representative}
If $M\subset U\subset\mathbb R^n$ is $\mathcal H^m$-measurable, $U$ is open, and $\mathcal H^m(M\cap K)<\infty$ for every compact $K\subset U$, there is a Borel set $M\subset B\subset U$ with $\mathcal H^m(B\setminus M)=0$. The Borel measure $\sigma(A)=\mathcal H^m(M\cap A)$ on $U$ is locally finite and Radon, and is unchanged if $M$ is replaced by $B$. The measures obtained by translation and dilation are likewise unchanged.
\end{lemma}
\begin{proof}
Choose increasing compact sets $K_j\subset U$ exhausting $U$. For each integer $j$, the set $M\cap K_j$ has finite Hausdorff measure. Borel regularity of Hausdorff measure, recalled in Definition \ref{def:hausdorff-measure}, provides a Borel superset $B_j$ with the same measure. Measurability and finite measure give $\mathcal H^m(B_j\setminus M)=0$ after subtracting $M\cap K_j$. The union $B=\bigcup_jB_j$ contains $M$ and differs from it by a countable union of null sets. Restrictions agree on every Borel test set, proving the measure identity and local finiteness. Intersect the Borel hull with $U$ if necessary. A locally finite Borel measure on a Euclidean open set is Radon. Pushforwards preserve equality of measures, giving the last assertion. This supplies a Borel representative when applying the Radon differentiation theorem; it does not identify topological boundaries of representatives.
\end{proof}

\begin{proposition}[The tangent measure of a continuously differentiable graph]\label{prop:c1-graph-tangent-measure}
Let $V\subset\mathbb R^m$ be open, let $u:V\to\mathbb R^{n-m}$ be $C^1$, and let $N=\{(z,u(z)):z\in V\}$. At $x=(a,u(a))$ with $a\in V$, put $T=\{(v,Du(a)v):v\in\mathbb R^m\}$. For every $\phi\in C_c(\mathbb R^n)$,
\[
 \lim_{r\downarrow0}r^{-m}\int_N\phi((y-x)/r)\,d\mathcal H^m(y)
 =\int_T\phi\,d\mathcal H^m.
\]
Only the portion of $N$ near $x$ enters this identity. This proves the smooth-patch limit used in \cite{Sim18}, p.~87, from the cited Euclidean area formula.
\end{proposition}
\begin{proof}
Write $J(z)=\sqrt{\det(I+Du(z)^*Du(z))}$. The weighted area formula, Proposition \ref{prop:weighted-area}, and the substitution $z=a+rv$ express the integral on the left as
\[
 \int_{(V-a)/r}\phi\left(v,\frac{u(a+rv)-u(a)}r\right)J(a+rv)\,dv.
\]
If the test is supported in $B_R(0)$, a contributing $v$ has $|v|\leq R$. For all sufficiently small $r$, this whole closed base ball lies in $(V-a)/r$. The segment formula and continuity of $Du$ give uniform convergence of the scaled height to $Du(a)v$ on that ball, and $J(a+rv)\to J(a)$ uniformly. The integrands are uniformly bounded on a fixed finite-volume ball. Dominated convergence gives
\[
 J(a)\int_{\mathbb R^m}\phi(v,Du(a)v)\,dv.
\]
The linear area formula for the injective graph map $v\mapsto(v,Du(a)v)$ identifies this with the integral on $T$. A sufficiently small ambient ball projects into a compact subset of $V$, so no global area bound on the whole graph is needed. The case $m=n$ has zero-dimensional target and $J=1$.
\end{proof}

\begin{corollary}[Unit density at an interior point of a graph patch]\label{cor:c1-patch-unit-density}
Under Proposition \ref{prop:c1-graph-tangent-measure},
\[
 \lim_{r\downarrow0}\frac{\mathcal H^m(N\cap B_r(x))}{\omega_m r^m}=1.
\]
The assertion applies at every interior point of a finite-area patch from Lemma \ref{lem:finite-area-c1-patches}.
\end{corollary}
\begin{proof}
The sphere $\partial B_1$ has zero $\mathcal H^m\llcorner T$ measure: in the $m$-plane it is an $(m-1)$-sphere. Sandwich its ball indicator between continuous cutoffs supported in slightly larger balls and equal to one on slightly smaller balls. Proposition \ref{prop:c1-graph-tangent-measure} then bounds the lower and upper limiting masses by $\omega_m(1-\delta)^m$ and $\omega_m(1+\delta)^m$. Let $\delta\downarrow0$. This proves the ball-mass limit without treating the discontinuous indicator as a continuous test.
\end{proof}

\begin{lemma}[Differentiation on the union of two measurable sets]\label{lem:union-measure-symmetric-difference}
Let $M,N\subset U\subset\mathbb R^n$ be Borel, with $U$ open, and have locally finite $\mathcal H^m$ measure in $U$, and put $\lambda=\mathcal H^m\llcorner(M\cup N)$. At $\mathcal H^m$-almost every $x\in M\cap N$,
\[
 \frac{\mathcal H^m(M\cap B_r(x))}{\lambda(B_r(x))}\to1,
 \qquad
 \frac{\mathcal H^m(N\cap B_r(x))}{\lambda(B_r(x))}\to1,
 \qquad
 \frac{\mathcal H^m((M\mathbin{\triangle}N)\cap B_r(x))}{\lambda(B_r(x))}\to0.
\]
If also $\mathcal H^m(N\cap B_r(x))/(\omega_mr^m)\to1$ at such a point, then the symmetric-difference numerator is $o(r^m)$ and $\lambda(B_r(x))/(\omega_mr^m)\to1$.
\end{lemma}
\begin{proof}
The union measure is locally finite and Radon. Theorem \ref{thm:radon-function-differentiation} applies to the bounded functions $\chi_M$ and $\chi_N$. At almost every point of their common part, both values are one and both averages tend to one. Their missing portions have relative measure tending to zero, and their sum is exactly the relative measure of $M\mathbin{\triangle}N$ within the union. This proves the first three limits. If the additional unit-density hypothesis holds, divide $\mathcal H^m(N\cap B_r)/(\omega_mr^m)$ by $\mathcal H^m(N\cap B_r)/\lambda(B_r)$ to obtain the limit for $\lambda$. Multiplication gives the $o(r^m)$ estimate. The union, rather than the sum of two restricted measures, avoids double-counting $M\cap N$.
\end{proof}

\begin{proposition}[Transfer of a tangent across an almost-full common part]\label{prop:rectifiable-patch-tangent-transfer}
Let $M\subset U$ be Borel with locally finite $\mathcal H^m$ measure in the open set $U$, and let $N\subset U$ be a finite-area $C^1$ graph patch. Then for $\mathcal H^m$-almost every $x\in M\cap N$, the set $M$ has the approximate tangent $T_xN$ in Definition \ref{def:approximate-tangent}. At these points,
\[
 \mathcal H^m((M\mathbin{\triangle}N)\cap B_r(x))=o(r^m).
\]
This accounts for the accumulating other pieces in the forward proof of \cite{Sim18}, Theorem 1.6.
\end{proposition}
\begin{proof}
Every point of $N$ is interior to its base domain. Corollary \ref{cor:c1-patch-unit-density} and Lemma \ref{lem:union-measure-symmetric-difference} give the stated small symmetric-difference mass outside an $\mathcal H^m$-null subset of $M\cap N$. Fix any such point. For a continuous test supported in $B_R(0)$, the absolute difference between the scaled $M$ and $N$ integrals is at most
\[
 \|\phi\|_\infty r^{-m}
 \mathcal H^m((M\mathbin{\triangle}N)\cap B_{Rr}(x))\longrightarrow0.
\]
The graph integral tends to the integral on $T_xN$ by Proposition \ref{prop:c1-graph-tangent-measure}. Therefore the $M$ integral has the same limit. The chosen exceptional set depends only on the symmetric-difference estimate, so one full-measure set works for every test and every fixed $R$, with no uncountable intersection of test-dependent null sets.
\end{proof}

\begin{theorem}[Rectifiable sets have approximate tangent planes]\label{thm:rectifiable-forward-tangents}
Let $M\subset U\subset\mathbb R^n$, with $U$ open, be $\mathcal H^m$-measurable, countably $m$-rectifiable and locally finite in $\mathcal H^m$ measure in $U$. Then the approximate tangent $T_xM$ exists at $\mathcal H^m$-almost every point of $M$. On each $C^1$ patch $N$ of a countable covering it equals $T_xN$ at almost every point of $M\cap N$. This is the forward direction of \cite{Sim18}, Chapter 3, Theorem 1.6 and Remark 1.5(2), pp.~87--88.
\end{theorem}
\begin{proof}
Replace $M$ by its Borel hull from Lemma \ref{lem:rectifiable-borel-representative}; this changes no restricted or scaled measure and preserves rectifiability up to a null set. Theorem \ref{thm:rectifiable-c1-decomposition} and Lemma \ref{lem:finite-area-c1-patches} cover the set up to a null subset by countably many Borel finite-area patches. Apply Proposition \ref{prop:rectifiable-patch-tangent-transfer} to each patch. The union of the exceptional subsets is still $\mathcal H^m$-null. At every other covered point the tangent exists and equals that of each patch containing it outside its exceptional subset. This proves both assertions. Local finiteness is used to make the union measure Radon; countable rectifiability alone does not supply it.
\end{proof}

\begin{corollary}[Density one on a locally finite rectifiable set]\label{cor:rectifiable-local-unit-density}
Under Theorem \ref{thm:rectifiable-forward-tangents},
\[
 \lim_{r\downarrow0}\frac{\mathcal H^m(M\cap B_r(x))}{\omega_mr^m}=1
 \quad\text{for }\mathcal H^m\text{-almost every }x\in M.
\]
The proof depends on Whitney extension, Euclidean area, and the Besicovitch input to Radon differentiation; it does not assume the full rectifiable-set tangent characterization.
\end{corollary}
\begin{proof}
At every point with a tangent, apply the continuous-cutoff argument of Corollary \ref{cor:c1-patch-unit-density} to the limit in Definition \ref{def:approximate-tangent}. The plane has unit normalized ball mass and does not charge the unit sphere. Theorem \ref{thm:rectifiable-forward-tangents} makes such points a full-measure set. This is the exact constant needed in the final perimeter measure-identification argument.
\end{proof}

\begin{corollary}[Tangent agreement on intersections]\label{cor:rectifiable-common-tangents}
If $M_1,M_2\subset U$ are measurable, countably $m$-rectifiable and locally finite in $\mathcal H^m$ measure in an open set $U\subset\mathbb R^n$, then
\[
 T_xM_1=T_xM_2
 \quad\text{for }\mathcal H^m\text{-almost every }x\in M_1\cap M_2.
\]
Their union is locally finite and countably rectifiable; it has the same tangent at almost every point of the intersection.
\end{corollary}
\begin{proof}
Use Borel representatives and let $\lambda$ be the Hausdorff measure restricted to their union. At almost every point of the intersection, both restricted measures have unit density by Corollary \ref{cor:rectifiable-local-unit-density}, and the symmetric-difference mass is $o(r^m)$ by Lemma \ref{lem:union-measure-symmetric-difference}. The scaled test integrals therefore have the same limits by the estimate in Proposition \ref{prop:rectifiable-patch-tangent-transfer}. The plane measures in these limits have supports equal to their planes, so equal limits imply equal planes. The mass of the union minus either set is bounded by the symmetric-difference mass, proving the final assertion by the same estimate. Intersecting the finitely many full-measure sets is legitimate.
\end{proof}
